\documentclass[10pt,a4paper]{article}
\usepackage[margin=21mm,headheight=14pt]{geometry}
\usepackage[T1]{fontenc}
\usepackage{lmodern,microtype,amsmath,amssymb,bm,booktabs,array,graphicx}
\usepackage{siunitx,xcolor,enumitem,fancyhdr,caption,float}
\usepackage{tikz}
\usetikzlibrary{arrows.meta,calc,decorations.markings}
\usepackage[hyphens]{url}
\usepackage{xurl}
\usepackage[colorlinks=true,linkcolor=blue!45!black,urlcolor=blue!45!black,citecolor=blue!45!black]{hyperref}
\usepackage[nameinlink,noabbrev]{cleveref}
\hypersetup{pdftitle={FFT2D5 validation in a uniform magnetic field},pdfauthor={OPALX numerical validation},pdfsubject={Thin-ellipsoid curved tracking benchmark and validation limits}}
\sisetup{detect-all}
\setlength{\parindent}{0pt}
\setlength{\parskip}{5pt plus 1pt}
\setlist{nosep,leftmargin=1.5em}
\captionsetup{font=small,labelfont=bf}
\pagestyle{fancy}\fancyhf{}
\fancyhead[L]{\small OPALX / FFT2D5 validation}
\fancyhead[R]{\small Uniform-field thin-bunch test}
\fancyfoot[C]{\thepage}
\newcommand{\SC}{space charge}
\newcommand{\T}{\mathsf T}
\newcommand{\dd}{\mathrm d}
\newcommand{\eps}{\epsilon_0}
\newcommand{\tabfile}[1]{\begin{center}\small\input{generated/#1.tex}\end{center}}
\newcommand{\reportfigure}[3]{\begin{figure}[H]\centering\includegraphics[width=\linewidth]{figures/#1.pdf}\caption{#2}\label{#3}\end{figure}}
\newcommand{\code}[1]{{\small\nolinkurl{#1}}}
\title{\vspace{-1.4em}\textbf{FFT2D5 validation in a uniform magnetic field}\\[0.4em]
\large A thin ellipsoid on a circular orbit: model, tracking tests and limits}
\author{OPALX numerical validation report}
\date{1 November 2026 \quad / \quad Working draft}

\begin{document}
\maketitle
\vspace{-1em}
\begin{abstract}
Actual OPALX tracking of a cold uniform-volume ellipsoid is compared with an
independent affine cold-fluid solution in a uniform magnetic field. The bunch
is long along its trajectory and initially thin perpendicular to the orbit.
With zero self-field and charges of 0.1, 0.3 and 1 pC, the fixed FFT2D5 solver
reproduces all three total RMS sizes within 1\% at every saved time step through
a half-turn. Independent mesh, time-step and particle-count controls support
this result. A radius similarity scan from 2 to 0.25 m also passes this size
criterion. However, the tangential space-charge-induced increase differs by
32\% at 1 pC. A completed 10 pC pilot makes this limitation more visible:
maximum total-size errors are 19.0\% radial, 5.65\% tangential and 2.59\%
vertical. Initial-field moments, covariance and centroid diagnostics identify
limits hidden by total RMS sizes. Strong-charge refinement and independent
3D controls are incomplete; no converged 10 pC accuracy claim is made.
\end{abstract}

\section{Geometry and coordinate convention}
The analytic problem uses physical coordinates $(X,Y,Z)$, a circular orbit in
the $X$--$Y$ plane, and a uniform applied field $\bm B=B_Z\bm e_Z$.
The identical problem is implemented in OPALX laboratory coordinates $(x,y,z)$
by a proper rotation and a translation:
\begin{equation}
 \begin{pmatrix}x+R\\y\\z\end{pmatrix}
 =T\begin{pmatrix}X\\Y\\Z\end{pmatrix},\qquad
 T=\begin{pmatrix}1&0&0\\0&0&1\\0&-1&0\end{pmatrix},\qquad
 T\bm B=\begin{pmatrix}0\\B_Z\\0\end{pmatrix},\quad \det T=1.
 \label{eq:coordinates}
\end{equation}
\textbf{Thus the executed OPALX test has an $x$--$z$ orbit, a thin $y$ extent,
and uniform $B_y$.} The label ``uniform $B_Z$'' refers to the physical analytic
coordinates. A global OPALX $B_z$ would not be perpendicular to this orbit.
In the results, $r$ denotes radial, $s$ tangential and $y$ vertical (OPALX).
Local reference-frame RMS values are distinct from global laboratory $x,z$
sizes, which change as the reference frame rotates.

\begin{center}\small
\begin{tabular}{ll}\toprule
Quantity & Baseline value\\\midrule
Species; centroid speed & Protons; $\beta_c=0.01$\\
Orbit radius; physical field & $R=1$ m; $B_Z=B_y^{\rm OPALX}=31.2973881$ mT\\
Initial semiaxes $(a_r,a_s,a_y)$ & $(5,50,0.5)$ mm\\
Initial RMS sizes $(\sigma_r,\sigma_s,\sigma_y)$ & $(2.23607,22.36068,0.223607)$ mm\\
Macroparticles; initial velocity & 98,304; rigid corotation, no random spread\\
Duration; steps & $\pi/\Omega=1.047922511~\mu$s; 128\\
FFT2D5 mesh; transverse box & $128\times256\times1024$; $24\times12$ mm\\
Longitudinal approximation & Existing OPEN model, fixed $g=6.36$\\\bottomrule
\end{tabular}
\end{center}
The initial centroid phase is 0.2 rad. Its orbit center in OPALX is
$(-R,0,0)$. Orbit coverage of $180^\circ$ must be distinguished from the
instantaneous bunch angular extent, initially about $5.7^\circ$ at $R=1$ m.

\clearpage
\section{Independent semi-analytic reference}
\Cref{fig:test-geometry} summarizes the orbit and initial bunch geometry used
by the independent model and OPALX. The local radial and tangential directions
rotate with the reference particle; the thin direction is normal to the orbit.
\begin{figure}[H]
 \centering
 \input{figures/test-geometry.tikz}
 \caption{Uniform-field test geometry in OPALX coordinates. Left: the reference
 centroid follows a clockwise half-turn when $+z$ points right and $+x$ up;
 $+B_y$ points out of the page and corresponds to physical $B_Z$.
 The bunch is shown at an illustrative intermediate phase, with outward
 $\bm e_r$ and forward $\bm e_s$. Right: two initial ellipsoid projections
 show the semiaxes, with separate magnifications. The orbital radius and bunch
 size are not drawn to a common scale.}
 \label{fig:test-geometry}
\end{figure}
The initial semiaxis hierarchy is $a_s:a_r:a_y=100:10:1$.
The centroid starts at phase $\theta_0=0.2$ rad and advances by $\pi$ rad;
the depicted intermediate bunch is not the initial phase. Rigid corotation
means that each particle's initial velocity depends on its position, with no
random velocity spread. The finite $y$ thickness allows a three-dimensional
uniform-volume ellipsoid reference; it is not an infinitesimal charged sheet.

\clearpage
\subsection{Cold uniform-volume ellipsoid}
For a random vector $\bm u$ uniformly distributed in the unit three-dimensional
ball, write the physical particle position and covariance as
\begin{equation}
 \bm r=\bm C+A\bm u,\qquad
 \rho=\frac{3Q}{4\pi\det A},\qquad
 \Sigma=\left\langle(\bm r-\bm C)(\bm r-\bm C)^\T\right\rangle
       =\frac15 AA^\T .\label{eq:affine}
\end{equation}
Here $A$ is dimensional (metres), $Q$ is the total bunch charge, and
$\det A>0$. This is a uniform \emph{volume} distribution with finite thickness;
its projection onto the orbit plane is not a uniform areal disk.

The electrostatic field inside the ellipsoid is exactly linear in this model:
\begin{equation}
 \bm E(\bm r)=\frac{\rho}{\eps}D(A)(\bm r-\bm C),\qquad
 AA^\T=U\operatorname{diag}(a_1^2,a_2^2,a_3^2)U^\T,\qquad
 D=U\operatorname{diag}(D_1,D_2,D_3)U^\T .\label{eq:field}
\end{equation}
The dimensionless depolarization factors are evaluated using Carlson's
symmetric elliptic integral $R_D$:
\begin{align}
 D_i&=\frac{a_1a_2a_3}{3}R_D(a_j^2,a_k^2,a_i^2),\qquad \sum_iD_i=1,\label{eq:carlson}\\
 D_i&=\frac{a_1a_2a_3}{2}\int_0^\infty
 \frac{\dd u}{(a_i^2+u)\sqrt{\prod_{\ell=1}^{3}(a_\ell^2+u)}} .
\end{align}
The equivalent integral provides an independent implementation check.
Affine cold-fluid ellipsoid dynamics in a uniform field are discussed by
Dubin~\cite{dubin}; the present benchmark has no electrostatic trap.

With $q=e$, $m=m_p$, and $\Omega=qB_Z/m$, the dimensional equations are
\begin{equation}
 J=\begin{pmatrix}0&1&0\\-1&0&0\\0&0&0\end{pmatrix},\qquad
 \ddot{\bm C}=\Omega J\dot{\bm C},\qquad
 \ddot A=\frac{q\rho}{m\eps}D(A)A+\Omega J\dot A .\label{eq:model}
\end{equation}
The sign gives clockwise physical rotation for positive protons in positive
$B_Z$. In OPALX, $TJT^\T$ gives acceleration proportional to $(-v_z,0,v_x)$.
The unrotated initial conditions are
\begin{align}
 A(0)&=\operatorname{diag}(5,50,0.5)~\mathrm{mm},&
 \dot A(0)&=\Omega J A(0),\\
 \bm C(0)&=(R,0,0)^\T,&
 \dot{\bm C}(0)&=(0,-\Omega R,0)^\T .
\end{align}
``Cold'' means zero random thermal spread about this deterministic velocity
field; particles do not all have the same velocity. At $Q=0$, the ellipsoid
rotates rigidly and its orbit-basis RMS sizes remain constant.

\subsection{Accuracy, assumptions and comparison axes}
The model assumes nonrelativistic, electrostatic, collisionless free-space
dynamics, with no magnetic self-field, conducting boundaries or radiation.
It is an exact affine solution of that continuum model while $A$ remains
nonsingular, integrated numerically by DOP853. At $\beta\simeq0.01$, omitted
relativistic and magnetic self-field terms are small; they are not declared
zero. OPALX uses relativistic momenta and has a separately checked zero-self-field
orbit. There is no fitted charge, field strength, time scale or $g$ in the oracle.

At each STAT time, the full covariance is rotated and projected onto the
\emph{actual OPALX reference axes}, derived from its recorded reference momentum.
With $R_0=\exp(\theta_0J)$ and $\theta_0=0.2$ rad, the OPALX laboratory
covariance is $\Sigma_{\rm OPALX}=T R_0AA^\T R_0^\T T^\T/5$.
\begin{equation}
 \sigma_i^{\rm model}(t)=\sqrt{\bm e_i^\T(t)\,\Sigma_{\rm OPALX}(t)\,\bm e_i(t)},
 \qquad i\in\{r,s,y\} .\label{eq:projection}
\end{equation}
This prevents reference-phase error from being mistaken for size error.
\tabfile{modelchecks}

\clearpage
\section{Numerics, implementation and test criteria}
\subsection{Fixed production implementation}
The tracking results use the native rectangular-cell-integrated two-dimensional
logarithmic Green function, $G(\bm r)=-\log|\bm r|/(2\pi)$, including its finite
self-cell average. Centered electric-field gradients use exterior potential
halos obtained from the doubled free-space convolution. This also avoids
artificial boundary self-forces. One FFT workspace is reused over slices;
exactly empty slices are skipped after clearing staged fields, while
zero-net-charge nonempty sources are still solved.

The particle/reference-path frame repair is included: the immutable laboratory
path and particles are transformed into a common solve frame and the particle
frame is restored after solving, including exceptional exits. The benchmark
uses the existing horizontal reference-path convention and longitudinal CIC.
No additional curvature correction or change of $g$ is part of these results.

The baseline reference path covers 0 to $4\pi/3$ rad (length about 4.18877 m).
The applied field is supplied by one constant Cartesian field map; its domain
contains the orbit and bunch. Open path endpoints must lie beyond the bunch,
including the deposition support. A segmented dipole pilot is not used as a
uniform-field reference because its sector seams failed that control.

\subsection{Metrics and acceptance}
The primary total-size metric is
\begin{equation}
 \epsilon_i=\max_{t\in\mathrm{STAT}}
 \left|\frac{\sigma_i^{\rm OPALX}(t)}{\sigma_i^{\rm model}(t)}-1\right| .
 \label{eq:error}
\end{equation}
Accepted baseline tracking requires $\epsilon_i<1\%$ in all three directions
at every promised step, complete time coverage, unchanged particle count and
correct assigned charge. Each independent baseline refinement must change
every RMS curve by at most 0.5\%, and must itself remain within 1\% of the model.
No tolerance was relaxed to obtain these results.

The space-charge response is a separate, more sensitive observable:
\begin{equation}
 \Delta\sigma_i(Q,t)=\sigma_i(Q,t)-\sigma_i(\mathrm{SC\ off},t),\qquad
 e_{\Delta,i}=\Delta\sigma_i^{\rm OPALX}-\Delta\sigma_i^{\rm model} .
 \label{eq:growth}
\end{equation}
Its relative error uses the model increase as denominator. Near zero increase,
absolute errors are reported instead of interpreting singular relative values.
Matched controls use the same source, field, radius and time step.

For HDF5 covariance checks, the standardized component error is
\begin{equation}
 e_{ij}^{\Sigma}=\frac{\Sigma_{ij}^{\rm OPALX}-\Sigma_{ij}^{\rm model}}
 {\sigma_i^{\rm model}\sigma_j^{\rm model}} .\label{eq:covmetric}
\end{equation}
Cross-covariances are measured only at retained particle snapshots; they cannot
be reconstructed from the three STAT RMS columns alone. Centroid response is
the charged centroid minus its matched self-field-off control. The independent
electrostatic model predicts zero charge-dependent centroid displacement.

\subsection{Space-charge-off control}
The NONE solver disables the self-field. A nominal finite assigned charge of
0.1 pC is retained so macroparticle charge-to-mass ratios are defined. Therefore
the control's STAT charge is \emph{not} zero; its model self-charge is zero.
Validators distinguish these two quantities.

\clearpage
\section{Completed FFT2D5 baseline intensity scan}
\tabfile{baseline}
\reportfigure{baseline-rms}{All-step total RMS sizes and signed differences from
the independent ellipsoid model. Model curves are explicitly identified;
size axes include zero. The maximum absolute error is taken over every STAT
sample, not only the final state.}{fig:baseline}
All four cases complete the half-turn and preserve all 98,304 particles and
their assigned charge. The largest baseline error is 0.90323\% in radial size
at 1 pC. This establishes the stated total-RMS criterion for this geometry.
It does not imply that the smaller collective increments are known to 1\%.

\clearpage
\section{Completed numerical refinements and 3D controls}
\subsection{Independent 1 pC FFT2D5 refinements}
\tabfile{convergence}
\reportfigure{baseline-refinement}{Independent changes to the accepted 1 pC run.
These are completed full-half-turn controls with the final production binary.
Agreement with the independent model and changes relative to the baseline are
different metrics.}{fig:convergence}
The largest refinement-induced RMS change is 0.31101\%. Every refined run
remains within 1\% of the model; the largest refined error is 0.98585\%.
These are separate one-parameter controls, rather than a simultaneous
asymptotic extrapolation of all discretizations.

\subsection{Completed 3D OPEN comparisons}
The existing three-dimensional OPEN solver provides an additional tracking
control for the same physical source, field and independent model. Its coarse
mesh is $32\times32\times64$; the refined 1 pC mesh is $64\times64\times128$.
The charge-off row below uses 512 steps; charged rows use 128.
\tabfile{open}
At 1 pC, halving the time step changes the largest size by 0.00831\%; doubling
particles changes it by 0.02515\%. Serial and two-rank OPEN curves differ by
less than $3.5\times10^{-13}$ relative. These distributed OPEN checks do not
establish distributed FFT2D5 support. Historical OPEN and development-run
binary identities remain in their original manifests; they are not relabeled
as runs of the final FFT2D5 binary.

\clearpage
\section{Tenfold-charge stress test: completed pilot}
The bunch charge is increased from 1 to 10 pC at unchanged source shape, radius,
field and half-turn duration. The model predicts final RMS sizes of
$(5.82549,24.81866,11.59273)$ mm in $(r,s,y)$. Its maximum vertical support
reaches 25.9 mm, so the transverse box is enlarged to $40\times64$ mm with a
$128\times512\times1024$ mesh. The baseline box would clip the expanded bunch.
\tabfile{strong}
\reportfigure{strong-charge-rms}{Completed 10 pC FFT2D5 pilot versus the independent
model. The stronger expansion exposes substantial radial and tangential
disagreement. The 1\% baseline acceptance is not achieved.}{fig:strong}
The maximum radial and tangential errors occur at the final time; the maximum
vertical error occurs earlier. All 128 STAT samples, particle count, assigned
charge and full duration are verified. All 16 retained particle snapshots have
zero particles outside the configured transverse box and zero in its outer
CIC half-cell; no claim is made about particle clearance at unsaved times.

This is a stress result, not a converged 10 pC accuracy measurement. The
fine-grid run and independent 10 pC OPEN control are incomplete. Moreover, the
model's final tangential-to-vertical support ratio is only about 2.14, so the
long-bunch approximation is stressed as the initially thin direction expands.
The discrepancy cannot yet be assigned solely to $g$ or declared a new code bug.

\clearpage
\section{Collective growth, covariance and centroid}
\subsection{Final increase relative to the matched control}
\textbf{1 pC:}\tabfile{growth1}
\textbf{10 pC pilot:}\tabfile{growth10}
The 1 pC tangential growth error is 32.36\%, despite a total tangential-size
error of only 0.3838\%. At 10 pC, the tangential increase is 2.45845 mm in the
model and 3.85746 mm in OPALX. Thus the stronger-charge test does increase the
space-charge signal and makes the model discrepancy much easier to see.
\reportfigure{collective-growth}{Collective RMS increases relative to actual
matched self-field-off controls. Dashed model and solid OPALX curves compare
the response directly; lower panels show signed absolute errors in that
increase. Size and growth axes include zero.}{fig:growth}

\clearpage
\subsection{Diagnostics beyond diagonal RMS sizes}
\tabfile{response}
The covariance column gives the final maximum absolute standardized component
error from \cref{eq:covmetric}; $\rho_{rs}$ is the radial--tangential correlation.
The centroid column is the maximum charged-minus-control displacement norm,
using the newer matched R1 self-field-off run. It is not the raw displacement
from a nonrelativistic reference circle.
\reportfigure{centroid-covariance}{Space-charge response beyond total RMS sizes. The
independent model has charge-independent centroid motion. Only retained HDF5
times are used for cross-covariances; unavailable times are not interpolated.}{fig:response}
For the 10 pC pilot, HDF5 and STAT RMS values agree within
$4.3\times10^{-15}$ relative at common snapshots. This checks the diagnostic
reconstruction but does not remove the physical/model disagreements.

\clearpage
\section{Initial self-fields and the longitudinal approximation}
The fixed initial source is evaluated with a single tiny-time-step field solve.
Fields are transformed into its principal axes and compared with
\cref{eq:field}. The all-particle electric-field moment ratio and particlewise
relative $L^2$ error are
\begin{equation}
 M_i=\frac{\sum_p \xi_{p,i} E_{p,i}^{\rm OPALX}}
 {\sum_p \xi_{p,i} E_{p,i}^{\rm exact}},\qquad
 L_i=\sqrt{\frac{\sum_p(E_{p,i}^{\rm OPALX}-E_{p,i}^{\rm exact})^2}
 {\sum_p(E_{p,i}^{\rm exact})^2}} .
\end{equation}
Here $\bm\xi_p$ is the initial displacement from the centroid. Subtracting the
applied field before evaluating a self-Lorentz-force diagnostic avoids counting
the bending magnet as a self-field.
\tabfile{fields}
Historical mesh: $64\times128\times1024$, binary prefix \code{415fff3b41df}.
Current mesh: $128\times256\times2048$, binary prefix \code{ba08fdd0b0bb}.
\textbf{This pair changes both mesh and executable; it is not a controlled
same-binary force-convergence study.} In particular, pointwise longitudinal
force convergence has not been established.
\reportfigure{initial-fields}{Initial field diagnostics. Historical and current
frozen solves are distinguished by provenance. A near-unity field moment does
not establish pointwise force accuracy.}{fig:fields}
\clearpage
\subsection{Shape-dependent longitudinal moment equivalent}
For the initial ellipsoid with radial, tangential and vertical semiaxes
$(a,b,c)$, matching the exact longitudinal field to the parabolic-line-density
convention gives $g_{\rm eq}=2b^2D_s/(ac)=5.223213$. The fixed configured
$g=6.36$ is 21.764\% larger. The current tangential moment is 21.281\% high,
consistent with an existing longitudinal-model difference already at the source.
For a tilted shape, the corresponding moment diagnostic is
\begin{equation}
 b_{\rm marg}=\sqrt{5\Sigma_{ss}},\qquad
 g_{\rm eq}(t)=\frac{2b_{\rm marg}^3}{\det A}
 \frac{(D\Sigma)_{ss}}{\Sigma_{ss}} .
\end{equation}
Here $D$ and $\Sigma$ are both expressed in the same ideal-centroid
radial--tangential--vertical basis before taking their $ss$ components.
It falls to 3.85049 at 1 pC and 1.90716 at 10 pC. The scalar does not reproduce
the full tilted-ellipsoid field tensor.
\reportfigure{longitudinal-model}{Longitudinal moment equivalent derived only
from the independent ellipsoid model, compared with the fixed configured
$g=6.36$. It is not fitted to OPALX and is not proposed as a scalar replacement
for the full tensor field.}{fig:gmodel}
The current initial mean radial
self-electric field is $-0.3432$ V/m rather than zero; a local circular-basis
decomposition attributes about 99.5\% of that mean to the longitudinal field
projection. This identifies a model diagnostic, not an implemented correction.

\clearpage
\section{Completed radius similarity scan}
Keep the source semiaxes, quiet sample and centroid $\beta$ fixed and choose
\begin{equation}
 B(R)=B_1\frac{1\,\mathrm m}{R},\quad
 Q(R)=Q_1\left(\frac{1\,\mathrm m}{R}\right)^2,\quad
 t(R)=t_1\frac{R}{1\,\mathrm m},\quad Q_1=1~\mathrm{pC}.
\end{equation}
With $\tau=\Omega t$ and a fixed normalization length $L$, the affine equation
depends on $\kappa=3qQ/(4\pi\eps m\Omega^2L^3)$. Since $Q/\Omega^2$ is
constant, the model predicts identical size-versus-angle curves. Each source's
rigid-corotation velocities, map and path are regenerated for its radius.

\tabfile{radius}
\reportfigure{radius-similarity}{Radius similarity test: all four charged cases and their
matched self-field-off controls complete the half-turn. The independent model
predicts coincident size-versus-angle curves. Total-size agreement remains
below 1\% in this scan.}{fig:radius}
The scan uses $N_x,N_y=128,256$ and $N_s=\operatorname{round}(1100R/\mathrm m)$,
preserving longitudinal physical spacing at about 4.09 mm. The reference path
is padded from $-0.15$ to 4.35 rad. All saved snapshots pass transverse-box and
CIC-edge clearance checks. The reported instantaneous bunch spans are about
$2.9^\circ,5.7^\circ,11.5^\circ,22.7^\circ$ in descending radius order;
these are span diagnostics, not general solver validity bounds.
\tabfile{radiusresponse}
The nearly constant tangential-growth bias points to an error already present
in the longitudinal approximation. The increasing centroid response shows that
curvature effects are not absent. The tight-radius fine-grid run is incomplete,
so the extreme-radius result has no completed independent refinement control.

\clearpage
\section{Software verification and automated tracking checks}
\begin{center}\small
\begin{tabular}{p{0.31\linewidth}p{0.61\linewidth}}\toprule
Test layer & Evidence and scope\\\midrule
Independent model & Sphere/oblate/integral, zero-charge rotation, centroid,
invariants, volume and tighter integration; \cref{eq:model}.\\
C++ unit tests / CTest & 131/131 registered entries pass with one OpenMP thread.
Kernel self-cell, anisotropy/scaling, quadrature, Gaussian force, halo/boundary,
small grids, neutral sources, stale fields, frames/containers/exceptions.\\
Four-thread numerical tests & 66 applicable FFT numerical tests pass; the
multiple-rank guard test is skipped in that single-rank invocation.\\
Two-rank FFT2D5 check & Explicit unsupported-MPI rejection passes on two ranks.
This is a configuration guard, not distributed tracking validation.\\
3D OPEN multi-rank control & Actual one-/two-rank tracking curves agree to
$3.5\times10^{-13}$ relative for the checked 1 pC case.\\
Native tracking regressions & ZeroSC and 1 pC cases registered locally in the
companion regression tree; independent all-step validators both pass.\\
Smoke-runner integration & ZeroSC executed through the new hook: normal
\code{.rt} checks and independent all-step validation pass. Seven runner
integration tests and meaningful validator corruption controls pass.\\
Remote CI & Workflow and companion cases prepared locally; unpublished and
not executed remotely. External CSCS runner remains on existing \code{.rt} checks.\\\bottomrule
\end{tabular}
\end{center}
The two native cases are \code{FFT2D5-UniformB-ZeroSC} and
\code{FFT2D5-UniformB-1pC}. Their validators check every one of 128 promised
STAT rows against independent covariance predictions, plus time coverage,
particle count, assigned charge and pinned source fixtures. Corrupting an
intermediate RMS while retaining the final value, truncating output, losing a
particle, changing charge/time or injecting nonfinite data must fail. The
three final-value \code{.rt} checks complement these checks; they do not replace them.

These are \emph{OPALX tracking regressions}. The solver C++ tests are unit tests.
The new full smoke hook has been exercised on ZeroSC; the 1 pC reference has
passed its independent validator, but has not been rerun through that hook.
GPU execution and multi-rank FFT2D5 tracking are unverified/unsupported.

\subsection{Incomplete cases excluded from acceptance}
\begin{center}\small
\begin{tabular}{lrrl}\toprule
Case & STAT rows & Duration covered & Intended check\\\midrule
10 pC fine mesh & 74/128 & 57.8125\% & $256\times1024\times1024$\\
10 pC 3D OPEN & 240/256 & 93.7500\% & Independent strong-charge control\\
$R=0.25$ m fine mesh & 61/128 & 47.6563\% & $256\times512\times550$\\\bottomrule
\end{tabular}
\end{center}
These are partial stored outputs, not pending claims of success. They have no
completion artifact and are excluded from completed-run comparisons. No new
simulation was started to assemble this report.

\clearpage
\section{Conclusions and next measurements}
The fixed implementation passes the specified 1\% \emph{total RMS size}
criterion for the thin, low-speed 0--1 pC baseline, with independent numerical
controls. The completed radius scan extends that observation across the tested
2--0.25 m geometries. This is meaningful curved tracking validation with a
rotating reference frame, not proof of a universal curved-space-charge model.

The charge-dependent observables remain substantially less accurate. At 1 pC,
tangential collective growth is 32\% high; the completed 10 pC pilot has 19\%
radial total-size error and 57\% tangential-growth error. The current initial
longitudinal field moment is about 21\% high. Covariance and a small nonzero
charge-dependent centroid shift give additional failures of the independent
electrostatic reference. These observations cannot be hidden by the successful
baseline RMS criterion.

The next measurements are to complete the 10 pC refinement and independent
OPEN control, complete the tight-radius refinement, and repeat a controlled
same-binary frozen-field study with adequate longitudinal sampling. Only then
can numerical uncertainty, the fixed longitudinal approximation and the
weakening long-bunch assumption be separated more decisively. No fitted $g$ or
additional curvature operator is introduced here.

\appendix
\section{Reproduction and provenance}
The main fixed implementation is saved in local commit \code{fd6ba8dc5} on
\code{658-fft2d5-keep-particles-and-reference-path-in-a-consistent-frame}.
Subsequent validation and reporting edits are local and uncommitted. The
archived final FFT2D5 tracking runs identify their executable by SHA-256
\begin{quote}\small\ttfamily
\path{ba08fdd0b0bb1e90a64ba0fd3763a8b6d044b128c0af9b9dfd215ad0ffac8fd4}.
\end{quote}
Historical archived provenance records the pre-commit parent plus exact source
hashes. It must not be read as proof that those files were unmodified at that
parent commit. The source content was subsequently saved in the local fix commit.
The coarse frozen solve uses a different historical binary, as labeled above.
These hashes identify the recorded runs, not whichever executable is currently
installed in a build directory.

This report's \code{data/provenance.json} lists source paths, byte counts and
SHA-256 hashes of copied tables, metadata, source code and logs. Figures and tables rebuild
from those local copies, without OPALX or network access. The large particle
HDF5 files and distributions remain in the original run tree; their manifests
and original analysis scripts preserve case provenance.

From this report directory, run \code{./build.sh}. It uses the OPALX Python
environment by default, regenerates vector plots and tables, and invokes
\code{latexmk -pdf}. Set \code{OPALX_REPORT_PYTHON} to another Python containing
NumPy, pandas and Matplotlib. The PDF is copied to
\code{output/pdf/fft2d5-uniform-field-validation.pdf} in the OPALX checkout.

The existing simulation generator and model remain one directory above this
report. For example, from that directory:
\begin{quote}\small\ttfamily\raggedright
python prepare\_opalx.py q1 --solver fft2d5 --charge-pc 1\newline
\hspace*{1em}--mesh 128 256 1024 --run\newline
python analyze\_stat.py --manifest opalx-runs/q1/manifest.json\newline
\hspace*{1em}--require-one-percent
\end{quote}
Use the configured OPALX environment and executable. An existing completed run
is preserved by the generator. Frozen fields and extended scans use their
archived drivers listed in the evidence catalog; rebuilding the report itself
does not launch any simulation.

\clearpage
\section{Evidence catalog}
Paths below are relative to the report's copied \code{data/} directory unless
marked as original run-tree paths. CSV tables and plot inputs are archived at
full stored precision; displayed results are rounded only for presentation.
\begin{center}\footnotesize
\begin{tabular}{>{\raggedright\arraybackslash}p{0.36\linewidth}>{\raggedright\arraybackslash}p{0.56\linewidth}}\toprule
Evidence & Files / original source\\\midrule
FFT2D5 accepted scan and refinements & \code{baseline-summary.csv},
\code{baseline-convergence.csv}, \code{baseline-acceptance.json}; parent
package \code{results/}.\\
3D OPEN controls & \code{open-summary.csv}, \code{open-refinement-summary.csv},
\code{open-convergence.json}; original \code{results/} and \code{convergence-analysis/}.\\
10 pC pilot & \code{strong-summary.csv}, \code{strong-manifest.json}; original
\code{strong-charge/pilot-analysis/}.\\
Collective response & \code{growth-summary.csv}, \code{centroid-summary.csv},
\code{covariance-snapshots.csv}; original \code{response-audit/results-q10-pilot/}.\\
Frozen fields / longitudinal equivalent & \code{field-summary.csv},
\code{g-evolution.csv}; original \code{response-audit/results-q10-pilot/}.\\
Radius scan & \code{radius-summary.csv}, \code{radius-all-case-summary.csv},
\code{analytic-geometry.csv}; original \code{radius-scan/results/}.\\
Model checks and code & \code{model-checks.json}, \code{strong-model-checks.json},
\code{ellipsoid_model.py}, \code{prepare_opalx.py}.\\
Unit/build verification & \code{test-logs/fft2d5-final-*.log}; original
\code{fft-goal/test-logs/}.\\
Native regression validation & \code{zero-regression-validation.json},
\code{q1-regression-validation.json}; companion \code{regression-tests-x}.\\
Partial cases & \code{incomplete-manifests/}; original strong-charge and radius
run directories, excluded from acceptance.\\\bottomrule
\end{tabular}
\end{center}
The original run-tree root is
\path{/Users/adelmann/git/OPALX-2D5/curved-study/uniform-bz}.
Historical coarse FFT2D5 development runs and failed field-map/seam pilots are
retained there. They diagnose development errors and are not substituted for
the final-binary acceptance scan. Straight KV tests and comparisons to a
different curvature operator are outside this report's uniform-field focus.

\begin{thebibliography}{9}
\bibitem{dubin}
D. H. E. Dubin, ``Equilibrium and dynamics of uniform density ellipsoidal
non-neutral plasmas,'' \emph{Physics of Fluids B} \textbf{5}, 295--324 (1993).
\href{https://nnp.ucsd.edu/pdf_files/PFB_5_295_1993.PDF}{Author-hosted full text}.
\end{thebibliography}
\end{document}
